% ---------------------------------------------
% 第2章 問題意識と研究目的
% ---------------------------------------------


\chapter[問題意識と研究目的]{問題意識と研究目的}
\label{chap:research_question}

第1章では，新電力市場の制度的背景として，電力自由化後の市場構造，
JEPX 価格形成，需給調整メカニズム，および制度変更が企業キャッシュフローに
与える影響を整理した。これらの制度的特徴は，新電力企業が他産業に比べて
外部ショックの影響を強く受ける構造的脆弱性を持つことを示している。

本章では，この制度的背景を踏まえ，新電力企業の倒産急増という現象を
従来モデルでは十分に説明できないという問題意識から出発し，
本研究が解明すべき倒産リスクの構造原理を明確化する。

% ---------------------------------------------
% 2.1 研究課題の所在
% ---------------------------------------------

\section{研究課題の所在}

2016 年の電力小売全面自由化以降，多数の新電力が市場に参入したが，
近年は倒産・撤退が急増している。

こうした倒産急増の背景には、財務指標では捉えきれない外部ショック型の倒産\footnote{
帝国データバンクおよび東京商工リサーチは、2021年以降の新電力倒産の特徴として
「財務指標が健全でも外部ショックにより短期間で資金繰りが破綻する」
外部ショック型倒産の存在を指摘している。
}
が存在する点を見落としてはならない。
日本市場では、帝国データバンクおよび東京商工リサーチの倒産報告\parencite{TDB2021PowerCrisis,TSR2021PowerCrisis}において、
以下のような事例が確認されている。


\begin{itemize}
    \item \textbf{日本の事例（A社）}\\  
    2021年1月のJEPX価格急騰により調達費用が短期間で急増し、直前期まで黒字であったにもかかわらず
    数週間で資金繰りが破綻した。財務指標が健全でも倒産に至る典型的な外部ショック型倒産である。

    \item \textbf{日本の事例（B社・C社）}\\  
    同じ外部ショックを受けたにもかかわらず、自治体からの資金支援を受けたB社は倒産を回避した一方、
    支援主体が存在しなかった\footnote{電気新聞および自治体議会資料によれば、自治体との包括契約に基づく短期資金支援の有無が倒産回避の可否を左右したと報告されている。}
C社は撤退に追い込まれた\parencite{DenkiShimbun2021Crisis,CityCouncil2021Support}。
    財務指標よりも「支援能力」の有無が企業存続を左右した事例である\footnote{新電力の倒産事例では、親会社・自治体・金融機関などの支援主体の存在が倒産回避の可否を決定づけると複数の倒産レポートで指摘されている。}
。
\end{itemize}

同様の構造は、制度構造の異なる英国市場でも観察される。
Ofgem の公式記録\parencite{Ofgem2021BulbSLR}によれば、燃料価格高騰期において
Bulb Energy は直前期の財務指標が急悪化していなかったにもかかわらず、
価格キャップ制度により調達費用の上昇を料金へ転嫁できず倒産に至った\footnote{Ofgem の公式記録では、価格キャップ制度の改定が小売事業者のキャッシュフローを圧迫し、倒産ラインを大きく押し上げた事例が複数報告されている。}
。
一方、Octopus Energy は親会社の資本支援と厚いヘッジ契約により同じショックを乗り切った\parencite{Ofgem2021OctopusTakeover,UKGov2021BulbSupport}。

これらの事例は、日本と英国という制度構造の異なる二つの市場において、
\textbf{財務指標が良好でも倒産する企業と生存する企業が分かれる構造的特徴が共通して存在する\footnote{日本と英国の新電力市場では、財務指標が直前期まで健全であっても外部ショックや制度コスト負担により倒産に至る事例が共通して観察される。}
}ことを示している。
すなわち、新電力企業の倒産は財務指標の線形的悪化ではなく、
外部ショック感応度、制度コスト負担、支援能力といった構造的要因の相互作用によって決定される可能性が高い。


特に 2021 年以降の JEPX 価格高騰，
需給調整金の急増，インバランス制度の厳格化といった外部ショックは，
新電力のキャッシュフローに直接的な負荷を与え，短期間で資金繰りが
悪化する事例が相次いだ。

TSR（東京商工リサーチ）や TDB（帝国データバンク）が公表する
倒産レポートでは，新電力の倒産には以下の特徴が共通して見られると指摘されている。

\begin{itemize}
    \item 財務指標が直前期まで良好であっても，外部ショックにより急速に資金繰りが悪化する
    \item 調達戦略やヘッジ比率の違いにより，外部ショックの影響度が大きく異なる
    \item 親会社・自治体・金融機関などの支援主体の有無が，倒産回避の可否を左右する
    \item 制度コスト（需給調整金・インバランス負担・担保制度等）の急増が倒産ラインを押し上げる\footnote{2021年以降の制度変更により、需給調整金やインバランス負担が急増し、平常時とは異なる倒産閾値が形成されたと指摘されている。}

\end{itemize}

これらの指摘は，新電力の倒産が単なる財務悪化の延長ではなく，
市場制度・外部環境・資本構造が複雑に絡み合う現象であることを示している。
すなわち，従来の財務指標のみでは倒産リスクを十分に説明できず，
外部ショック感応度，制度コスト負担，支援能力といった構造的要因を
考慮する必要がある。

本章では，官報に掲載された決算公告に基づき，
代表的な 5 社（F-Power，ウエスト電力，シナジアパワー，
エフエネ，洸陽電機〔現シン・エナジー〕）の財務推移を比較し，
倒産企業・倒産寸前企業・生き残り企業の構造的特徴を整理する。
これら 5 社の決算公告は付録に掲載している。

\begin{table}[H]
\centering
\caption{市場観察対象 5 社の特徴（定性的整理）}
\begin{tabular}{p{3cm}p{2cm}p{8cm}}
\toprule
企業名 & 状態 & 主な特徴（定性的） \\
\midrule
F-Power & 倒産 &
急成長後に外部ショックの影響を強く受け，調達費用の急騰に耐えられず破綻。財務指標は直前期まで良好であったが，市場価格変動に対する脆弱性が顕在化した。 \\

ウエスト電力 & 倒産 &
利益率の低下と負債負担の増大により財務基盤が弱体化。外部ショック以前から収益性が低く，制度コスト増加が資金繰りを圧迫し破綻に至った。 \\

シナジアパワー & 倒産 &
短期負債への依存度が極めて高く，長期資金調達力が乏しい構造。売上規模は拡大したが，調達費用の急騰により逆ザヤが拡大し，短期間で債務超過が深刻化した。 \\

エフエネ & 倒産寸前→再生 &
数期連続の大幅債務超過に陥ったが，外部からの資本支援により再生。財務指標だけでは倒産に至らず，支援主体の存在が企業存続を左右した典型例。 \\

洸陽電機（現シン・エナジー） & 生き残り &
資産規模が大きく，長期借入を含む安定的な資金調達が可能。外部ショック期においても黒字を維持し，財務基盤が強固で市場変動に耐性を有した。 \\
\bottomrule
\end{tabular}
\end{table}

各社の財務推移を観察すると，それぞれに固有の事情が存在する。
F-Power は市場価格変動の影響を強く受けて破綻し，
ウエスト電力は収益性の低下と負債負担の増大が倒産の主因となった。
シナジアパワーは短期負債依存が極めて高く，調達費用の急騰により
短期間で債務超過が拡大した。エフエネは深刻な債務超過に陥ったものの，
外部からの資本支援により再生している。洸陽電機（現シン・エナジー）は，
資本力と長期資金調達力を背景に外部ショック期を乗り切った生き残り企業である。

一方で，これらの企業には共通する特徴も見られる。
第一に，倒産企業の多くは，財務指標が直前期まで良好であっても，
外部ショックにより急速に資金繰りが悪化している。
第二に，市場価格変動や制度コストの増加といった外部環境の変化が，
企業の財務構造に即時に波及している。
第三に，親会社・金融機関・スポンサーなどの支援主体の有無が，
企業の存続可能性に大きな影響を与えている。

これらの観察結果は，TSR や TDB が倒産レポートで指摘する
新電力倒産の特徴と整合的であり，新電力の倒産が単なる財務悪化ではなく，
外部ショック・制度コスト・支援能力といった構造的要因により
規定されている可能性を示唆する。本章の市場観察は，
これらの構造的要因を明らかにするための基礎的分析として位置づけられる。

加えて、新電力企業の倒産は、日本市場に特有の制度的脆弱性だけでなく、
自由化が先行した英国市場でも同様の構造的特徴が確認されている。
英国では、価格キャップ制度\footnote{
英国の価格キャップ制度（Default Tariff Price Cap）は，
2018 年の *Domestic Gas and Electricity (Tariff Cap) Act* に基づき導入された制度であり，
規制当局 Ofgem が標準的家庭向け料金（default tariff）に対して上限価格を設定する仕組みである。
上限価格は，燃料費，ネットワークコスト，環境負担金，運営費などの原価要素を基に
半年ごと（2022 年以降は四半期ごと）に改定される。

この制度は「消費者保護」を目的としているが，
卸電力価格が急騰した局面では，小売事業者が調達コストを販売価格に転嫁できず，
逆ザヤが発生し，キャッシュフローが急速に悪化する構造的問題が指摘されている。
2021～2022 年の燃料価格高騰期には，価格キャップが実勢価格に追随できず，
多くの新電力が赤字を抱え倒産に至った。

なお，Ofgem の公表資料によれば，
2021 年だけで 28 社以上の小売事業者が倒産しており，
その多くが「価格キャップによる収益圧迫」を主要因として挙げている。
}，
SLR（Supplier of Last Resort）制度，卸市場の変動性などが企業のキャッシュフローに直接波及し，
外部ショック型の倒産が多数発生している\footnote{
英国の SLR（Supplier of Last Resort）制度は、倒産した小売事業者の顧客を他社が引き継ぐ際に，
引き受け企業へ追加負担が発生する制度であり，燃料価格高騰期にはこの負担が急増し，
倒産ラインを押し上げる要因となったと指摘されている。
}。
本研究では、日本市場で観察された倒産メカニズムが、制度構造の異なる英国市場でも再現されるのかを検証することで、
企業の倒産リスクが制度差を超えて普遍的な構造原理に基づくものであるかを明らかにする。

% ---------------------------------------------
% 2.2 従来モデルの限界
% ---------------------------------------------

\section{従来モデルの限界}

倒産研究の伝統的アプローチは，財務比率を説明変数とする線形モデルである。
代表的なものとして，Altman（1968）の Z スコアモデル\footnote{
Altman（1968）の Z スコアモデルは，多変量判別分析（MDA: Multiple Discriminant Analysis）
を用いて倒産企業と非倒産企業を識別する統計モデルである。
製造業を中心とした上場企業の財務データを基に構築され，
運転資本比率，利益率，市場価値ベースの自己資本比率など複数の財務指標を
線形結合して「倒産スコア」を算出する点に特徴がある。
モデルは倒産予測の古典的手法として広く普及し，信用リスク管理の基礎的枠組みとして位置づけられている。
}や，
Ohlson（1980）の O スコアモデル\footnote{
Ohlson（1980）の O スコアモデルは，ロジットモデルを用いて倒産確率を直接推定する手法であり，
Altman の判別分析モデルに比べて統計的仮定が緩やかで，より一般的な産業に適用可能とされた。
モデルは財務比率だけでなく，企業規模，財務構造，収益性，資本構成など
複数の非線形要因を含む説明変数を用いる点に特徴があり，
倒産確率を連続的に推定できることから，実務の信用リスク評価でも広く利用されてきた。
}が挙げられる。
これらのモデルは，流動比率，自己資本比率，利益率などの財務指標を
加算的に組み合わせることで倒産確率を推定するものであり，
製造業や小売業など多くの産業で一定の予測精度を示してきた。


従来モデルが広く受け入れられてきた背景には，
倒産が「財務悪化の延長線上で徐々に進行する」という前提がある。
すなわち，収益性の低下や負債の増加といった財務指標の悪化が，
時間をかけて倒産リスクを押し上げるという線形的なプロセスである。
近年では機械学習を用いた非線形モデルも登場しているが，
依然として財務指標に依拠する点では同様である。

しかし，2.1節で整理したように，
新電力企業の倒産急増は，従来モデルでは十分に説明できない特徴を有している。

第一に，\textbf{外部ショックの影響が極めて大きい}点が挙げられる。
JEPX 価格の急騰や需給調整金の増加といった外部環境の変化は，
企業のキャッシュフローに即時に波及し，短期間で資金繰りを悪化させる。
従来モデルは市場価格変動や制度コストといった外生的ショックを
明示的に扱わないため，外部環境の急変による倒産を捉えにくい。

第二に，\textbf{倒産閾値が固定されていない}という特徴がある。
新電力市場では，制度変更や需給逼迫に伴うキャッシュアウトの急増により，
倒産ラインそのものが平常時と大きく異なる。
従来モデルは倒産閾値を一定と仮定するため，
制度コストの急増による倒産リスクの跳ね上がりを反映できない。

第三に，\textbf{支援能力の有無が企業存続を左右する}点を適切に評価する必要がある。
親会社・自治体・金融機関などの支援主体の存在は，
倒産回避の可否を決定づけるが，従来モデルは外生的な支援能力を
説明変数として扱わない。
そのため，財務指標が悪化していても支援により再生する企業や，
逆に財務指標が良好でも支援主体が不在で倒産に至る企業を説明できない。

第四に，\textbf{複数要因が同時に悪化する複雑性}がある。
新電力の倒産は，市場価格，制度コスト，資本構造，事業運営の脆弱性など，
複数の要因が相互に関連しながら進行する。
しかし従来モデルは，説明変数を加算的に扱うため，
要因間の相互作用や掛け算的な悪化プロセスを捉えにくい。

以上より，新電力企業の倒産リスクは，従来の財務比率モデルでは捉えきれない
構造的特徴を持つことが明らかである。本研究では，これらの限界を踏まえ，
外部ショック，制度コスト，支援能力，財務構造の相互作用を統合的に捉える
新たな倒産リスクモデルの構築を目指す。

なお，従来モデルの限界は日本市場だけでなく，英国市場でも同様に指摘されている。

英国規制当局 Ofgem は，2021～2022 年の新電力倒産について次のように述べている \parencite{Ofgem2022SupplierFailures}。
\begin{quote}
“the unprecedented rise in wholesale gas prices exposed suppliers with insufficient hedging and weak capital positions.”
\end{quote}
（日本語訳）  
「前例のない卸ガス価格の高騰は，十分なヘッジを持たず資本構造の弱い事業者の財務ストレスを一気に増幅させた。」

また，Supplier of Last Resort（SoLR）発動に伴う顧客引き継ぎコストについても，
Ofgem は次のように指摘している \parencite{Ofgem2022SupplierFailures}。
\begin{quote}
“the SoLR process imposed significant costs on the industry, accelerating financial distress.”
\end{quote}
（日本語訳）  
「SoLR プロセスは業界全体に多額のコストを課し，財務的な困難を加速させた。」

英国会計検査院（NAO）も，同期間の倒産要因について次のように報告している \parencite{NAO2022SupplierFailures}。
\begin{quote}
“the cost of protecting customers through the SoLR process reached £2.7bn.”
\end{quote}
（日本語訳）  
「SoLR による顧客保護コストは 27 億ポンドに達した。」

さらに，英国ビジネス・エネルギー省（BEIS）は，価格キャップ制度について次のように述べている \parencite{BEIS2022RetailStrategy}。
\begin{quote}
“the current price cap design can amplify supplier financial stress during periods of wholesale volatility.”
\end{quote}
（日本語訳）  
「現在の価格キャップ制度は，卸価格の変動期に小売事業者の財務ストレスを増幅し得る。」

学術研究においても同様の指摘が存在し，
Hall and Roelich（2022）は次のように述べている \parencite{HallRoelich2022EnergyPolicy}。
\begin{quote}
“supplier failures in 2021 were driven primarily by regulatory design, not by traditional financial deterioration.”
\end{quote}
（日本語訳）  
「2021 年の新電力倒産は，従来型の財務悪化ではなく，規制設計が主因であった。」

\vspace{1\baselineskip}

これらの国際的な観察結果は，倒産リスクが財務指標の単純な悪化ではなく，
制度的ショックと資本構造の相互作用によって非線形的に決定されるという
本研究の問題意識を強く裏付ける。


% ---------------------------------------------
% 2.3 問題意識
% ---------------------------------------------

\section{問題意識}

前節で整理したように，新電力企業の倒産は，
外部ショック，制度コスト，資本構造，支援能力といった
複数の要因が相互に作用することで発生しており，
従来の財務比率モデルでは十分に説明できない特徴を有している。
この観察結果を踏まえると，本研究が取り組むべき中心的な問題意識は，
次の三点に集約される。

\begin{enumerate}
    \item \textbf{財務指標が良好でも倒産するという違和感}\par
          官報公告に基づく市場観察では，倒産企業の多くが
          直前期まで黒字や正の純資産を維持しており，
          従来モデルが前提とする「財務悪化の延長としての倒産」と整合しない。
          この点は，TSR や TDB の倒産レポートが指摘する
          「外部ショック型倒産」の特徴とも一致する。

    \item \textbf{外部ショックが即時に経営へ波及するという特殊性}\par
          新電力市場は，JEPX 価格や制度コストの変動が
          企業キャッシュフローに直接的かつ短期間で影響する
          “高感応型市場”である。
          このような外部環境の急変は，従来モデルが前提とする
          緩やかな財務悪化プロセスとは大きく異なる。

    \item \textbf{複数の構造的要因が同時に作用するという複雑性}\par
          市場価格，制度コスト，資本構造，事業運営の脆弱性，
          さらには支援主体の有無など，多様な要因が相互に関連しながら
          倒産リスクを形成している。
          このような掛け算的・相互作用的な倒産プロセスは，
          加算的な財務比率モデルでは捉えにくい。
\end{enumerate}

以上の点から，新電力企業の倒産リスクは，
単一の財務指標ではなく，複数の構造的要因の組み合わせによって決まると考えられる。
本研究では，この問題意識を出発点として，
新電力市場に特有の倒産メカニズムを理論的に再構成することを目指す。

以上の観察結果は、新電力企業の倒産リスクが、制度差や市場構造の違いを超えて、企業の財務構造に内在する普遍的な脆弱性によって規定されている可能性を示唆する。すなわち、倒産リスクの本質は、国や制度に依存するものではなく、外部ショック感応度・資本構造・支援能力といった企業固有の構造的要因の組み合わせによって決まるという点にある。本研究は、この普遍的構造原理を理論的に再構成し、後章で日本企業と英国企業のデータを用いてその妥当性を検証する。

% ============================================================
% 2.4 研究動機（筆者の企業経験）
% ============================================================


\section{研究動機（筆者の企業経験）}

本研究の背景には、筆者自身の企業経験に基づく問題意識がある。
筆者は四十余年にわたり一企業に勤務し、その間に三度、会社存続にかかわる重大な危機に直面した。

一度目は、1985 年 9 月のプラザ合意を契機として 1986～1987 年に急激な円高が進行し、
輸出依存型の事業構造が大きな収益悪化に直面した時期である。
二度目は、競合他社との技術革新競争において優位性を確保できず、
製品競争力の低下が事業継続を脅かした時期である。
三度目は、2008 年のリーマンショック後の金融環境悪化により、
メインバンクからの与信が得られず資金繰りが逼迫した時期である。

いずれの局面においても、企業は多面的な努力により倒産を回避し、
筆者自身も一社での職業人生を全うすることができた。
しかし、もしこれらの危機が倒産に至っていたならば、
家族への負担、地域社会への影響、取引先や従業員を含むステークホルダーへの損失は計り知れない。
この経験は、企業倒産が単なる経済現象ではなく、
広範な社会的影響を伴う重大事象であることを強く意識させる契機となった。

新電力市場は制度的に参入障壁が低く、競争原理が導入されている一方で、
淘汰される企業も多い。企業努力のみならず、制度設計そのものが倒産リスクにどのように影響するのかを検証することは、
市場の持続可能性を考えるうえで重要である。
筆者の経験に照らせば、倒産回避のメカニズムを学術的に明らかにすることは、
企業経営のみならず社会的安定に資する可能性がある。

以上の問題意識から、本研究では新電力企業の倒産リスクを段階的に捉えるモデルを構築し、
日本市場に加えて英国市場を対象とした実証分析を行い、
制度的特徴と倒産リスクの関係を深掘りすることを目的とする。



% ---------------------------------------------
% 2.5 研究目的
% ---------------------------------------------

\section{研究目的}

前節までの検討から，新電力企業の倒産リスクは，
外部ショック，制度コスト，資本構造，支援能力といった
複数の要因が相互に作用することで形成されており，
従来の財務比率モデルでは十分に説明できないことが明らかとなった。
この問題意識を踏まえ，本研究の目的は，
新電力企業の倒産リスクがどのような構造原理によって決定されるのかを，
市場要因・財務要因・資本構成要因を統合した枠組みによって明らかにすることである。

具体的には，以下の点を解明する。

\begin{itemize}
    \item 外部ショック（市場価格変動・制度コスト）が倒産リスクに与える影響の大きさ
    \item 調達戦略や顧客構成といった事業運営上の脆弱性が倒産過程に果たす役割
    \item 親会社・自治体・金融機関などによる支援能力が企業存続に与える決定的影響
    \item 制度コストの急増が倒産閾値をどのように押し上げるかというメカニズム
\end{itemize}

これらの要因を統合的に捉えることで，
新電力企業の倒産リスクが単一の財務指標ではなく，
複数の構造的要因の組み合わせによって決まることを理論的に示す。
本研究は，新電力市場に特有の倒産メカニズムを構造的に再構成し，
従来モデルでは捉えきれなかった倒産リスクの本質を明らかにすることを目指す。

さらに、本研究では、日本市場で構築した倒産リスクモデルを英国市場にも適用し、
制度構造の異なる市場においても同様の構造原理が成立するかを検証する。
英国市場は自由化の歴史が長く、価格変動リスクや制度コストの急増が新電力企業の倒産に直結する点で日本市場と類似しており、
モデルの外挿可能性とロバスト性を確認する対象として適切である。
これにより、倒産リスクの構造原理が国境を越えて普遍的に成立するかを明らかにする。


% ---------------------------------------------
% 2.6 本研究の位置づけ
% ---------------------------------------------

\section{本研究の位置づけ}

本研究は，「新電力会社の資本構成と倒産リスク－財務要因と市場要因の包括的分析－」
という題目が示すように，新電力企業の倒産リスクを，
財務要因・市場要因・資本構成要因を統合した視点から包括的に分析する点に独自性がある。

従来の倒産研究は，財務比率を説明変数とする線形モデルが主流であり，
倒産リスクを財務指標の加算的な組み合わせとして捉えてきた。
しかし 2.1 節および 2.2 節で整理したように，
新電力企業の倒産は，外部ショック，制度コスト，資本構造，支援能力といった
複数の要因が相互に作用することで発生しており，
財務指標のみを前提とする従来モデルでは十分に説明できない。

本研究は，この従来モデルの限界を踏まえ，
市場価格変動や制度コストといった市場要因，
負債構成や資本調達力といった財務・資本構成要因，
さらに親会社・自治体・金融機関による支援能力といった外生的要因を
統合的に扱うことで，新電力市場に特有の倒産メカニズムを理論的に再構成する。

以上の問題意識と研究目的を踏まえると，
新電力企業の倒産リスクは，単一要因ではなく，
複数の構造的要因が相互に関連しながら形成される複雑な現象であると考えられる。
次章では，これらの構造的要因を整理し，
新電力市場に特有の倒産メカニズムを理論的に構造化したうえで，
本研究が検証すべき仮説（H1〜H4）を提示する。

また、本研究は、日本市場と英国市場という制度構造の異なる二つの市場を対象とし、倒産リスクの構造原理が国や制度を超えて普遍的に成立するかを検証する点に独自性がある。日英の制度差は倒産プロセスの表面的な違いを生むものの、企業の財務悪化が外部ショック・制度コスト・資本構造の相互作用によって進行するという根本構造は共通している可能性が高い。本研究は、この普遍的構造原理を理論的に示し、実証分析によりその妥当性を検証する。
